\subsection{Characteristic exponent of a field. Perfect fields}

Let K be a field. By the characteristic exponent of K we understand the integer equal to 1 if K is of characteristic 0, and equal to the characteristic of K when this is non-zero.

PROPOSITION 4. --- Let K be a field of characteristic exponent q. For every integer \(f \geqslant 0\) the mapping \(x \mapsto x^{q^{f}}\) is an isomorphism of K on one of its subfields (denoted by \(K^{q^{f}}\) ).

This follows from Th. 2 when \(q \neq 1\) and is trivial when \(q = 1\) .

Likewise one can extend Prop. 2 to the case where A is a field of characteristic exponent q, the case q = 1 being trivial.

DEFINITION 4. --- A field K of characteristic exponent q is said to be perfect if we have \(K^{q} = K\) . When \(K^{q} \neq K\) , K is called imperfect.

By this definition a field is perfect if it is of characteristic 0, or if it is a perfect ring of characteristic \(p \neq 0\) in the sense of Def. 2. If \(\mathbf{K}\) is a field of characteristic \(p \neq 0\) and \((\hat{\mathbf{K}}, u)\) is a perfect closure of \(\mathbf{K}\) , then \(\hat{\mathbf{K}}\) is a field by Prop. 3 (V, p.~6) and \(u\) is an isomorphism of \(\mathbf{K}\) onto a subfield of \(\hat{\mathbf{K}}\) . Frequently one identifies \(\mathbf{K}\) with its image under \(u\) in \(\hat{\mathbf{K}}\) , so that we have \(\hat{\mathbf{K}} = \mathbf{K}^{p^{-\infty}}\) (Prop. 3).

Let K be a field of characteristic 0; the characteristic exponent of K is then 1. By convention the notation \(x^{q^{-f}}\) and \(S^{q^{-f}}\) is taken to mean x and S respectively (for an element x of K and a subset S of K). In particular we put \(K^{q^{-\infty}} = K\) and we agree to take the perfect closure of K to be K.

PROPOSITION 5. --- If K is a field of characteristic 0, or is finite, * or algebraically closed *, it is perfect. In particular every prime field is perfect.

Suppose that K has characteristic \(p \neq 0\) . If K is finite, then the subfield \(K^{p}\) of K has the same cardinal as K, whence \(K^{p} = K\) . * If K is algebraically closed, the polynomial \(X^{p} - a\) has a root x in K for each \(a \in K\) (V, p.~20, Def. 1) whence \(x^{p} = a\) and so \(K^{p} = K\) . * Finally a prime field is of characteristic 0 or finite.

Let \(K_{0}\) be a field of characteristic \(p \neq 0\) and \(\mathbf{K} = \mathbf{K}_{0}(\mathbf{X})\) the field of rational fractions in an indeterminate X over \(K_{0}\) . Then K is imperfect, for there exists no element \(u(\mathbf{X})/v(\mathbf{X})\) of K (u, v polynomials in \(K_{0}[\mathbf{X}]\) ) such that \((u(\mathbf{X})/v(\mathbf{X}))^{p} = \mathbf{X}\) . This may be seen by writing this relation in the form \(u(\mathbf{X})^{p} = \mathbf{X}v(\mathbf{X})^{p}\) and comparing the degrees of the two sides.

